% v4 Methods 3.1: the near-body Chang lift-element head (MC-10).
% The full construction (the two equations, the discretisation, the QC gate and
% the proxy control) moved to appendix A alongside its own result (APA-07) and
% its figure (fig:phi), so that all three supervision heads are described at the
% same depth in the body. The \subsubsection heading went with it: it was the
% only head with one. Every constant from outputs/session35/mc_provenance.md
% MC-10; QC values through the macros pipeline.

\label{sec:methods_chang}

\rb{K}{S31-01}{The near-body head at the same depth as the other two: what the observable is, why it is not the wake target, and where the construction lives. Full construction moved to appendix A.}The near-body head targets a physically constructed observable rather than a raw field, because not all vorticity contributes equally to the lift. The force-element decomposition of \citet{chang1992}, in the form popularised for vortex-dominated flows by \citet{menon2021}, weights the vorticity at each point by an auxiliary potential that measures how strongly motion there projects onto the lift direction, so the resulting lift-element density is large where the flow is actually producing load and small where it is not. We supervise that density inside a feathered band within $0.3c$ of the surface, since the cache window truncates the domain and only the band-restricted field tracks the lift, and we compress it to $80$ dimensions by the identical construction used for the wake target above. The two field observables are therefore byte-matched in capacity and share one network class, and differ only in where they look: the wake head at the shed vortex layout, the near-body head at the vorticity that is generating the load. The construction, its discretisation and its pre-committed quality gate are given in appendix~\ref{app:nearbody}.\re
